\chapter*{초록}
\thispagestyle{plain}
SO-ARM101 로봇팔로 다섯 개 블록을 180초 이내에 지정 영역으로 옮기고, 후속 미션에서는 블록을 쌓아 5초간 유지하는 시스템을 개발하였다. 초기에는 SmolVLA와 ACT를 시험했으나, 최종 시스템에는 고정 카메라 기반의 고전 컴퓨터 비전, 평면 좌표 변환과 역기구학을 적용하였다. 색상·형상 기반 대상 선택, 중립 손목 방향을 고려한 파지 후보 생성, 그리퍼 위치·부하를 이용한 파지 확인, 재시도와 동적 슬롯 배치를 유한상태기계로 연결하였다.

초기 180초 제한 시험에서는 파지 후보 14회를 실행하였다. 센서 판정은 파지 4회, 빈손 9회였으며 1회는 판정 전에 중단되었다. 네 번의 슬롯 해제 후에도 마지막 블록이 영역 밖에 남아 미션을 완료하지 못했다. 별도의 단일 블록 시험에서는 파지와 운반에 성공했지만, 블록 위치를 바꾼 뒤에는 빈손 파지가 발생했다. 이후 그리퍼 회전 재시도와 하강 민감도 조정을 추가한 재시험에서는 28회 모두 180초 안에 다섯 개 블록을 배치하였다. 블록 최초 접근 217회 가운데 148회가 첫 시도에 성공하여 최초 파지 성공률은 68.2\%였다.

중간 아홉 점 캘리브레이션의 적합 RMS는 8.18\,mm, 최대 LOO 오차는 26.64\,mm였다. 로봇팔 여섯 축의 나사를 다시 조여 백래시를 줄이고, Z축 높이와 손목 회전각을 일정하게 유지하여 대응점을 재수집한 결과 최종 RMS는 5.15\,mm, 최대 LOO 오차는 13.53\,mm로 감소하였다. RMS는 개발 목표인 5\,mm 미만에 근소하게 미달했지만 중간 결과보다 37.1\% 줄었다. 한 블록의 첫 PICK에 51.5초가 걸린 사례도 있어 재시도에 드는 시간을 함께 고려해야 했다. 적재 시험 전에는 ArUco 마커와 손목 카메라를 장착했으며, 이후 그리퍼 회전 재시도와 하강 민감도 조정을 추가하였다. 3단 적재는 시험한 범위에서 모두 성공했다. 4단 적재는 추가 조정 후 1회 시도에 성공하여 해제 뒤 5초 이상 유지했지만, 5단 적재는 아직 시도하지 못했다. CV+IK 실행을 기록하는 데이터셋 자동 수집 기능도 구현하였다. 첫 실행에서 다섯 색상별로 1개씩 총 5개 에피소드(2{,}111프레임)를 186초 만에 폐기 없이 저장하였다. 다만 제어 주기는 목표보다 26\% 느렸으며, 1라운드의 결과로 수집 성능을 일반화할 수는 없다.

\medskip
\noindent\textbf{주요어:} SO-ARM101, 컴퓨터 비전, 역기구학, 유한상태기계, 파지 검증
